\begin{abstract}
We present a spectroscopy method for driven two-level systems based on a root-Lorentzian envelope with a $\pi$-phase inversion at its peak. For driven-spectroscopy pulses shorter than the lifetime $T_1$, the echo-like symmetry suppresses resonant excitation while preserving the off-resonant response. Measurements on a superconducting circuit compare ordinary and echo pulses through $20~\mathrm{MHz}$ nominal Rabi frequency, and pulse-shaping variants through $\LatestThreeDisplayMaxRabi~\mathrm{MHz}$ complex-envelope peak drive. Combining derivative removal by adiabatic gate (DRAG) with a polynomial correction reduces the median normalized linewidth $\mathrm{FWHM}/\Gamma_{T_2^*}$ from $\LatestDragWidthNormalized$ to $\LatestPolyWidthNormalized$, where $\Gamma_{T_2^*}=1/(\pi T_2^*)$ and $T_2^*=\PaperTTwoStar~\mu\mathrm{s}$, and increases median dip contrast from \LatestDragContrast{} to \LatestPolyContrast{} on common accepted amplitudes. The corrected center variation is $\LatestPolyCenterRms~\mathrm{kHz}$ RMS about its mean. These measured linewidths remain above the Ramsey coherence reference. Together with numerical studies of the linewidth--contrast tradeoff, the results support shaped echo pulses as a practical route to amplitude-robust spectroscopy beyond weak drive.
\end{abstract}
